\documentclass[11pt,a4paper]{article}
\usepackage[utf8]{inputenc}
\usepackage[T1]{fontenc}
\usepackage{lmodern}
\usepackage[french]{babel}
\usepackage{amsmath,amssymb,amsthm}
\usepackage[margin=2.4cm]{geometry}
\usepackage{xcolor}
\usepackage{tcolorbox}
\usepackage{booktabs}

\setlength{\parindent}{0pt}
\setlength{\parskip}{4pt}
\numberwithin{equation}{section}

\newcommand{\sign}{\operatorname{sign}}
\newcommand{\hth}{\hat{\theta}}
\newcommand{\hw}{\hat{\omega}}

\newtcolorbox{idee}{colback=blue!4,colframe=blue!50!black,title=Idée,fonttitle=\bfseries}
\newtcolorbox{resultat}{colback=green!5,colframe=green!40!black,title=Résultat,fonttitle=\bfseries}
\newtcolorbox{attention}{colback=red!4,colframe=red!50!black,title=Attention,fonttitle=\bfseries}

\title{\textbf{Pourquoi normaliser le détecteur de phase ?}\\[4pt]
\large Justification de $e \triangleq \varepsilon/\hat V_m$}
\date{}

\begin{document}
\maketitle

%=====================================================================
\section{Ce n'est pas une hypothèse, c'est un choix de conception}
%=====================================================================

L'expression
\begin{equation}
e \triangleq \frac{\varepsilon}{\hat V_m}
\end{equation}
n'est ni une hypothèse sur le réseau, ni un résultat que l'on démontre.
C'est une \textbf{définition} : on \emph{décide} de diviser le signal du détecteur par l'amplitude estimée.
Le symbole $\triangleq$ signifie justement \og est défini comme\fg.

\begin{idee}
On fait ce choix parce que le signal brut $\varepsilon$ pose un problème, que l'on explique ci-dessous.
La normalisation est la solution à ce problème.
\end{idee}

%=====================================================================
\section{Le problème : $\varepsilon$ dépend de l'amplitude du réseau}
%=====================================================================

\subsection{Deux informations mélangées}
Le détecteur de phase fournit
\begin{equation}
\varepsilon = V_m\sin\psi,\qquad \psi = \theta-\hth .
\end{equation}
Ce signal contient deux informations de nature différente :
\begin{itemize}
  \item $\sin\psi$ : \textbf{l'erreur de phase}, c'est ce que la PLL doit corriger ;
  \item $V_m$ : \textbf{l'amplitude du réseau}, qui n'a rien à voir avec la phase.
\end{itemize}

\subsection{Conséquence sur la boucle fermée}
Supposons que l'on utilise $\varepsilon$ directement dans la loi Super-Twisting :
\begin{equation}
\hw = k_1|\varepsilon|^{1/2}\sign(\varepsilon) + z + \omega_{ff},\qquad \dot z = k_2\sign(\varepsilon).
\end{equation}

\textbf{Étape 1.} On remplace $\varepsilon = V_m\sin\psi$ dans le terme en racine :
\begin{equation}
|\varepsilon|^{1/2} = |V_m\sin\psi|^{1/2} = V_m^{1/2}\,|\sin\psi|^{1/2},
\end{equation}
car $V_m>0$, donc $|V_m\sin\psi| = V_m|\sin\psi|$, et $(ab)^{1/2} = a^{1/2}b^{1/2}$ pour $a,b\ge0$.

\textbf{Étape 2.} Le signe ne change pas, car $V_m>0$ :
\begin{equation}
\sign(\varepsilon) = \sign(V_m\sin\psi) = \sign(V_m)\,\sign(\sin\psi) = \sign(\sin\psi).
\end{equation}

\textbf{Étape 3.} En reprenant le calcul de $\dot x_1 = \omega-\hw$ avec $x_1=\psi$ et $x_2=(\omega-\omega_{ff})-z$ :
\begin{equation}
\boxed{\dot x_1 = -\underbrace{k_1\sqrt{V_m}}_{\text{gain effectif}}\,|\sin x_1|^{1/2}\sign(x_1) + x_2,
\qquad \dot x_2 = -k_2\sign(x_1)+\dot\omega .}
\end{equation}

\begin{attention}
Le gain réellement appliqué n'est plus $k_1$ mais $k_1\sqrt{V_m}$. Il dépend de la tension du réseau,
que l'on ne contrôle pas et qui peut varier fortement.
\end{attention}

%=====================================================================
\section{Deux exemples concrets}
%=====================================================================

\subsection{Exemple 1 : réseau 230 V}
Pour un réseau 230 V efficaces, l'amplitude crête vaut
\begin{equation}
V_m = 230\sqrt2 \approx 325\ \text{V},\qquad \sqrt{V_m}\approx 18 .
\end{equation}
Le gain effectif vaut donc
\begin{equation}
k_1\sqrt{V_m}\approx 18\,k_1 .
\end{equation}
Avec $k_1=30$ (valeur utilisée dans la preuve), le gain réellement appliqué serait d'environ $540$.
Les gains calculés dans l'analyse de stabilité ne correspondent plus du tout au système réel.

\subsection{Exemple 2 : creux de tension}
Lors d'un défaut, la tension peut chuter à $20\,\%$ de sa valeur nominale : $V_m\to0.2\,V_m$.
Le gain effectif est alors multiplié par
\begin{equation}
\sqrt{0.2}\approx0.45 ,
\end{equation}
c'est-à-dire \textbf{divisé par plus de 2}.

La condition de stabilité $k_2>K(k_1,L)$ a été établie pour une valeur fixe de $k_1$.
Si le gain effectif change, cette condition peut ne plus être vérifiée,
et ce, précisément pendant un défaut, quand la PLL doit rester stable.

%=====================================================================
\section{La solution : diviser par l'amplitude}
%=====================================================================

\subsection{L'amplitude se calcule à partir des mesures}
\textbf{Étape 1.} On calcule la somme des carrés des composantes de Clarke :
\begin{equation}
V_\alpha^2+V_\beta^2 = V_m^2\cos^2\theta + V_m^2\sin^2\theta = V_m^2\left(\cos^2\theta+\sin^2\theta\right) = V_m^2 .
\end{equation}
\textbf{Étape 2.} On prend la racine carrée ($V_m>0$) :
\begin{equation}
\boxed{\hat V_m = \sqrt{V_\alpha^2+V_\beta^2} = V_m .}
\end{equation}
Cette quantité se calcule à chaque instant, sans connaître $\theta$ : c'est simplement la longueur du
vecteur $(V_\alpha,V_\beta)$.

\subsection{Définition de l'erreur normalisée}
\textbf{Étape 1.} On divise $\varepsilon$ par $\hat V_m$ :
\begin{equation}
e \triangleq \frac{\varepsilon}{\hat V_m} = \frac{-V_\alpha\sin\hth+V_\beta\cos\hth}{\sqrt{V_\alpha^2+V_\beta^2}} .
\end{equation}
\textbf{Étape 2.} On remplace $\varepsilon=V_m\sin\psi$ et $\hat V_m = V_m$ :
\begin{equation}
e = \frac{V_m\sin\psi}{V_m} = \sin\psi .
\end{equation}

\begin{resultat}
\begin{equation}
e = \sin\psi .
\end{equation}
L'amplitude $V_m$ a disparu. Il ne reste que l'information de phase.
\end{resultat}

\subsection{Conséquence sur la boucle fermée}
Avec $e$ à la place de $\varepsilon$, le même calcul qu'à la section 2 donne
\begin{equation}
\dot x_1 = -k_1|\sin x_1|^{1/2}\sign(x_1) + x_2,\qquad \dot x_2 = -k_2\sign(x_1)+\dot\omega .
\end{equation}
Le gain effectif est exactement $k_1$, quelle que soit la tension du réseau.

%=====================================================================
\section{Ce que l'on gagne}
%=====================================================================

\begin{center}
\renewcommand{\arraystretch}{1.35}
\begin{tabular}{lll}
\toprule
 & \textbf{Sans normalisation ($\varepsilon$)} & \textbf{Avec normalisation ($e$)} \\
\midrule
Expression & $V_m\sin\psi$ (en volts) & $\sin\psi$ (sans unité) \\
Plage de valeurs & $[-V_m,\,V_m]$ & toujours $[-1,\,1]$ \\
Gain effectif & $k_1\sqrt{V_m}$ (variable) & $k_1$ (constant) \\
Creux de tension & dynamique modifiée & aucun effet \\
Preuve de stabilité & dépend de $V_m$ & indépendante de $V_m$ \\
\bottomrule
\end{tabular}
\end{center}

\begin{idee}
Le terme $k_2\sign(\cdot)$ ne profite pas de la normalisation, et n'en a pas besoin :
comme $V_m>0$, $\sign(\varepsilon)=\sign(e)$. Seul le terme en $k_1$ est concerné.
\end{idee}

%=====================================================================
\section{Précautions pratiques}
%=====================================================================

\subsection{Division par zéro}
Si la tension s'effondre ($\hat V_m\approx0$, court-circuit franc), la division devient impossible.
On protège le calcul avec un seuil $V_{\min}>0$ :
\begin{equation}
e = \frac{\varepsilon}{\max\left(\hat V_m,\,V_{\min}\right)} .
\end{equation}
Dans ce cas extrême, le détecteur ne contient de toute façon plus d'information de phase.

\subsection{Réseau déséquilibré}
Si le réseau est déséquilibré, $\hat V_m$ n'est plus exactement constant : il oscille à la fréquence $2\omega$.
La normalisation reste utile, mais elle est imparfaite. C'est une raison supplémentaire d'ajouter un
pré-filtrage (DSOGI, MAF) avant le détecteur.

\subsection{Une technique classique}
La normalisation par l'amplitude est une pratique courante dans la littérature sur les PLL.
Elle permet de régler la dynamique de la boucle indépendamment du niveau de tension.
Ce n'est pas une astuce propre à cette preuve.

%=====================================================================
\section{Résumé}
%=====================================================================

\begin{resultat}
On n'a rien supposé. On a \textbf{choisi} de diviser $\varepsilon$ par une amplitude que l'on sait calculer
à partir des mesures, $\hat V_m=\sqrt{V_\alpha^2+V_\beta^2}$. Ainsi :
\begin{itemize}
  \item l'erreur normalisée vaut exactement $e=\sin\psi$ ;
  \item le gain de la boucle est exactement $k_1$ ;
  \item la PLL et sa preuve de stabilité ne dépendent plus de la tension du réseau.
\end{itemize}
\end{resultat}

\end{document}
